
A well-documented failure mode of empirical risk minimization (ERM) is the collapse of worst-group accuracy (WGA) on datasets with spurious correlations. On Waterbirds \citep{Sagawa2019Distributionally} --- where 95\% of training samples associate waterbirds with water backgrounds and landbirds with land backgrounds --- a ResNet-50 trained with ERM achieves high average accuracy while performing poorly on the minority groups that do not carry the spurious feature (56 waterbird-on-land samples and 184 landbird-on-water samples out of 4,795 total training samples).

The dominant supervised remedy, Group Distributionally Robust Optimization (GroupDRO) \citep{Sagawa2019Distributionally}, explicitly minimizes worst-group loss and achieves strong WGA, but requires per-sample group annotations. This requirement restricts practical deployment, since obtaining group labels demands domain expertise and is costly at scale.

Existing label-free alternatives use coarse proxy signals. Just Train Twice (JTT) \citep{Liu2021Just} uses binary misclassification by an ERM model as a minority proxy, upweighting the error set in a second training stage. Learning from Failure (LfF) \citep{Nam2020Learning} uses the relative cross-entropy loss between a biased and a debiased network to upweight hard samples. Neither provides a mechanistic account grounded in the gradient structure of ERM training.

This work investigates the per-sample normalized last-layer gradient norm $\tilde{g}_i$ as a minority group proxy. The mechanistic hypothesis is that the outer-product decomposition of the fully connected layer gradient, $\|\nabla _W \ell _{i}\|_F = \|p_i - y_{i}\| \cdot \|h(x_{i})\|$, directly exposes the prediction residual $\|p_i - y_{i}\|$, and that BatchNorm feature equalization renders $\|h(x_{i})\|$ approximately uniform across groups. Under this account, minority samples that cannot exploit the spurious shortcut maintain elevated prediction residuals throughout early ERM training, producing persistently elevated $\tilde{g}_i$ values relative to majority samples.

The reported experiments address four questions: (RQ1) Does $\tilde{g}$ achieve AUC > 0.70 for minority group prediction at T\_id = 5? (RQ2) What is the minority/majority ratio at T\_id $\in$ {1, 3, 5, 10}? (RQ3) Does the ratio persist at T\_id = 10? (RQ4) Does BatchNorm equalize feature norms (h\_norm\_std\_ratio < 0.5)?

\textbf{Experimental scope.} Only the existence experiment (H-E1) was executed. The proposed full pipeline --- Stage 2 last-layer retraining to measure WGA --- was not run. No WGA results are reported. The findings reported here concern the proxy signal quality only.

\textbf{Contributions:}

\begin{enumerate}
\item Empirical measurement of $\tilde{g}$ as a minority proxy signal: AUC = 0.914 and ratio = 8.8x at T\_id = 5, without group annotations, on Waterbirds.
\item Mechanistic confirmation via outer-product decomposition and BatchNorm feature equalization (h\_norm\_std\_ratio $\approx$ 0.10).
\item An efficient computation procedure using a forward hook on the FC layer, requiring no additional backward passes.
\item Identification of a criterion design error: class balance deviation is an inappropriate metric for minority-focused selection on imbalanced datasets; minority recall in the selected subset is the appropriate criterion.
\end{enumerate}

